\chapter*{\centering{Tóm tắt đồ án}}
\addcontentsline{toc}{chapter}{Tóm tắt đồ án}

Hệ thống băng tải phân loại màu sắc là một ứng dụng phổ biến, 
tuy nhiên vẫn còn tồn tại thách thức lớn trong việc đảm bảo tốc độ 
vận hành ổn định và chính xác.\\
\indent Đây là yếu tố then chốt để định vị 
và phân loại vật thể thành công, đặc biệt khi hệ thống phải đối phó
với nhiễu động từ môi trường. Để giải quyết vấn đề này, 
nhóm nghiên cứu đã xây dựng một hệ thống nhúng dựa trên 
vi xử lý ESP32, tích hợp cảm biến màu sắc, động cơ servo và 
cảm biến vật cản.\\
\indent Giải pháp cốt lõi được áp dụng là thuật toán 
điều khiển PID (Proportional-Integral-Derivative), cho phép hệ 
thống tự động tính toán và điều chỉnh công suất động cơ để băng 
tải đạt được và duy trì tốc độ mong muốn với sai số tối thiểu, 
qua đó tối ưu hóa đáng kể hiệu suất và độ tin cậy của quá trình 
phân loại. Việc ứng dụng PID để điều chỉnh tốc độ băng truyền đã 
giúp hệ thống phân loại đạt độ chính xác.